% Einheit 1 von 4 – Länge, Masse, Geld (bank/einheiten/e1.jsonl, zusammenbau v0.1)
%% TODO Zweigzeile Teil 1: was hier gelernt wird (Satz in Schülersprache aus dem Einheitstitel) – Einheit 1
\einheitenkopf[e1]{Einheit 1 von 4 · Länge, Masse, Geld}
\zweigzeile{ab Kl. 5 · P10}
\setcounter{aufgabe}{13}

%% TODO Titel: Ich-kann-Satz fehlt in der Bank (Platzhalter: Kettenname) – Nr. 14
%% TODO Anweisung: ein Satz über der ersten Teilaufgabe fehlt in der Bank – Nr. 14
\begin{aufgabe}{Passt die Angabe?}
\begin{teile}
\teil Wie hoch ist eine Zimmertür? Kreuze an.\\ \kreuz{$2$ m}\\ \kreuz{$9$ m}\\ \kreuz{$25$ cm}
\end{teile}
\end{aufgabe}

%% TODO \verfahren{…}: Verfahrensüberschrift in Sachsprache fehlt in der Bank (2.3 b)
%% TODO Titel: Ich-kann-Satz fehlt in der Bank (Platzhalter: Kettenname) – Nr. 15
%% TODO Anweisung: ein Satz über der ersten Teilaufgabe fehlt in der Bank – Nr. 15
\begin{aufgabe}{Länge, Masse und Geld umrechnen}
%% TODO \rechenplatz{n}: Schrittzahl des Grundfalls fehlt in der Bank (nur bei fester Schreibform, 2.3 b)
\begin{teile}
\teil Welche Einheit ist größer? Kreuze an.\\ \kreuz{Meter}\\ \kreuz{Millimeter}
\teil $7$ m – wie viel dm? \leerfeld[dm]
\teil $30$ mm – wie viel cm? \leerfeld[cm]
\teil $7\,500$ mm – wie viel m? \leerfeld[m]
\teil $3$ km – wie viel m? \leerfeld[m]
\teil $5$ t – wie viel kg? \leerfeld[kg]
\teil $3{,}45$ € – wie viel ct? \leerfeld ct
\teil $3$ m $45$ cm – wie viel cm? \leerfeld[cm]
\teil $2{,}35$ m – wie viel cm? \leerfeld[cm]
\teil Setze $<$, $>$ oder $=$ ein: $0{,}04$ m \leerfeld $40$ cm \hfill (P10 2019 OS)
\teil Ein Fernseher hat eine Bildschirmdiagonale von $55$ Zoll; $1$ Zoll $\approx 2{,}54$ cm. Weise nach, dass die Diagonale etwa $1{,}40$ m lang ist. \hfill (P10 2016 OS)
\end{teile}
\end{aufgabe}

%% TODO \verfahren{…}: Verfahrensüberschrift in Sachsprache fehlt in der Bank (2.3 b)
%% TODO Titel: Ich-kann-Satz fehlt in der Bank (Platzhalter: Kettenname) – Nr. 16
%% TODO Anweisung: ein Satz über der ersten Teilaufgabe fehlt in der Bank – Nr. 16
\begin{aufgabe}{Maßstab im Koordinatensystem, Sek II}
%% TODO \rechenplatz{n}: Schrittzahl des Grundfalls fehlt in der Bank (nur bei fester Schreibform, 2.3 b)
\begin{teile}
\teil Eine Längeneinheit entspricht $30$ cm. Gesucht ist die Höhe eines Regals. Unterstreiche den Maßstabssatz. Was wird umgerechnet? Kreuze an.\\ \kreuz{Länge}\\ \kreuz{Fläche}\\ \kreuz{Volumen}
\teil Eine Leiste ist $6$ Längeneinheiten lang; eine Längeneinheit entspricht $25$ cm. Wirkliche Länge? \leerfeld[cm]
\teil Die Kante eines Bauteils läuft von $A(1|2)$ bis $B(1|7)$; eine Längeneinheit entspricht $4$ cm. Wie lang ist die Kante? \leerfeld[cm]
\teil Ein Ball fliegt auf der Bahn $f(x) = -0{,}25x^2 + 2x + 1{,}5$ ($x$: waagerechter Abstand vom Abwurf in m, $f(x)$: Höhe in m). Ein Zaun steht $3$ m vom Abwurf entfernt. In welcher Höhe fliegt der Ball über den Zaun? \leerfeld[m]
\teil Eine Vase entsteht durch Drehung von $f(x) = 0{,}25x^2 + 1$ um die $x$-Achse; die Öffnung liegt bei $x = 2$, eine Längeneinheit entspricht $3$ cm. Durchmesser der Öffnung? \leerfeld[cm]
\teil Der Bogen einer Toreinfahrt wird durch $f(x) = 0{,}6x^4 - 3x^2 + 2{,}4$ zwischen den Nullstellen $x = -1$ und $x = 1$ beschrieben; eine Längeneinheit entspricht $1{,}6$ m. Im Tor stehen $7$ senkrechte Stangen im Abstand von $40$ cm, die mittlere in der Tormitte. Berechne die Länge der vorletzten Stange rechts und die Länge einer waagerechten Stange in $2{,}4$ m Höhe. \hfill (FHR 2025) \\ vorletzte Stange: \leerfeld[m] , waagerecht: \leerfeld[m]
\end{teile}
\end{aufgabe}

%% TODO Titel: Ich-kann-Satz fehlt in der Bank (Platzhalter: Kettenname) – Nr. 17
%% TODO Anweisung: ein Satz über der ersten Teilaufgabe fehlt in der Bank – Nr. 17
\begin{aufgabe}{Einheiten der Länge, der Masse und des Geldes der Größe nach ordnen}
\begin{teile}
\teil Ordne, mit der kleinsten beginnend: m, mm, km, cm, dm \leerfeld ; \leerfeld ; \leerfeld ; \leerfeld ; \leerfeld
\end{teile}
\end{aufgabe}

%% TODO Titel: Ich-kann-Satz fehlt in der Bank (Platzhalter: Kettenname) – Nr. 18
%% TODO Anweisung: ein Satz über der ersten Teilaufgabe fehlt in der Bank – Nr. 18
\begin{aufgabe}{Repräsentanten zuordnen (welche Angabe passt zu welchem Gegenstand)}
\begin{teile}
\teil Ordne zu: Heft, Fahrrad, Fußballfeld – $105$ m; $30$ cm; $1{,}8$ m Heft \leerfeld , Fahrrad \leerfeld , Fußballfeld \leerfeld
\end{teile}
\end{aufgabe}

%% TODO Titel: Ich-kann-Satz fehlt in der Bank (Platzhalter: Kettenname) – Nr. 19
%% TODO Anweisung: ein Satz über der ersten Teilaufgabe fehlt in der Bank – Nr. 19
\begin{aufgabe}{drei bis vier Größen ordnen}
\begin{teile}
\teil Ordne, mit der kleinsten beginnend: $1{,}2$ m; $95$ cm; $1\,150$ mm \leerfeld ; \leerfeld ; \leerfeld
\end{teile}
\end{aufgabe}

%% TODO Titel: Ich-kann-Satz fehlt in der Bank (Platzhalter: Kettenname) – Nr. 20
%% TODO Anweisung: ein Satz über der ersten Teilaufgabe fehlt in der Bank – Nr. 20
\begin{aufgabe}{Umrechnen mit gegebenem Faktor bei nichtmetrischen Einheiten (Fuß, Meile, Zoll; P10-Form)}
\begin{teile}
\teil Ein Monitor misst $32$ Zoll in der Diagonale; $1$ Zoll $\approx 2{,}54$ cm. Wie viel cm? Runde auf eine Stelle. \leerfeld[cm]
\end{teile}
\end{aufgabe}

%% TODO Titel: Ich-kann-Satz fehlt in der Bank (Platzhalter: Kettenname) – Nr. 21
%% TODO Anweisung: ein Satz über der ersten Teilaufgabe fehlt in der Bank – Nr. 21
\begin{aufgabe}{Größen mit gleicher Einheit addieren und subtrahieren}
\begin{teile}
\teil $3{,}45$ m $+ 1{,}8$ m? \leerfeld[m]
\end{teile}
\end{aufgabe}

%% TODO Titel: Ich-kann-Satz fehlt in der Bank (Platzhalter: Kettenname) – Nr. 22
%% TODO Anweisung: ein Satz über der ersten Teilaufgabe fehlt in der Bank – Nr. 22
\begin{aufgabe}{Länge, Masse und Geld umrechnen – Fehler finden, Begründen, Anwendung, Darstellungswechsel}
\begin{teile}
\teil Tom schreibt: $450$ cm $= 45\,000$ m. Wo ist der Fehler?
\teil Eine Längeneinheit entspricht $0{,}4$ cm. Ein Bauteil ist $7$ Längeneinheiten hoch. Jan schreibt: Das Bauteil ist $7$ cm hoch. Wo ist der Fehler?
\teil Begründe: Warum ist $0{,}09$ m weniger als $90$ cm?
\teil Anna sagt: „Eine Längeneinheit ist immer ein Zentimeter.“ Begründe, warum das nicht stimmt.
\teil Ein Regal soll in eine Nische von $1{,}35$ m Breite. Die Böden sind $128$ cm lang, dazu kommen zwei Seitenwände von je $2$ cm. Passt das Regal hinein?
\teil Schreibe die Länge aus der Stellenwerttafel als Kommazahl in m und als Angabe in m und cm.

\sachtabelle{cccc}{m & dm & cm & mm}{2 & 0 & 6 & 0}
\leerfeld[m] ; \leerfeld[m] \leerfeld[cm]
\end{teile}
\end{aufgabe}
